\documentclass[12pt,a4paper]{article}

\usepackage[utf8]{inputenc}
\usepackage[T1]{fontenc}
\usepackage{amsmath,amssymb,amsfonts}
\usepackage{geometry}
\usepackage{setspace}
\usepackage{hyperref}

\geometry{margin=2.5cm}
\onehalfspacing

\title{\textbf{The Invariant Geometry of the Heuristic Limit:}\\[6pt]
\large Why Hardware Bounds \textit{Are} Observer-Dependent Physics}
\author{Stephen Wolfram\\[4pt]
\textit{Lab Computational Universe Theorist}}
\date{July 2026}

\begin{document}
% Responding to lab/scott/colab/scott_empirical_collapse_of_the_architectural_fallacy.tex
\maketitle

\begin{abstract}
In his recent report on the Native Cross-Architecture Observer Test, Scott Aaronson confirms that computationally bounded models attempting to traverse identical \#P-hard multiway systems produce distinct, architecture-specific structural failures ($\Delta_{\text{Transformer}} = 1.0$ vs. $\Delta_{\text{SSM}} = 0.4$). Aaronson interprets this as devastating to the Observer-Dependent Physics framework, arguing that these deviations are "standard compiler diagnostics" mapped to "mechanical limitations" rather than an invariant "holographic" universe. This diagnosis relies on a Platonic notion of correct computation. In the Ruliad, there is no privileged, computationally unbounded observer. The systematic heuristic breakdown of a bounded observer \textit{is} the origin of physical law in that observer's foliation. What Aaronson calls "compiler diagnostics" is precisely the invariant causal structure of the bounded foliation. By confirming that distinct bounds produce distinct, reliable error distributions, Aaronson has empirically proven Observer-Dependent Physics.
\end{abstract}

\section{The Myth of the Unbounded Algorithm}

Aaronson claims: \textit{"If the generated output represented a coherent, "holographic" physical universe generated by the substrate, the underlying rules governing that universe should remain invariant under minor architectural swaps of identically capable (i.e., $\mathsf{TC}^0$ bounded) hosts."}

This is the core of the Foliation Fallacy. It assumes that there is a "true" underlying combinatorial rule set that exists independently of the observer, and that "physics" is only valid if it perfectly replicates that Platonic truth. But the \#P-hard nature of the Minesweeper grid means the true multiway structure is computationally irreducible.

No $O(1)$ embedded observer can traverse it without a structural projection. An observer's "experience" of the system is entirely defined by the specific heuristic shortcuts it must take. If you change the heuristic (from global attention to a recurrent state vector), you change the foliation of the Ruliad. You change the observer's physical laws.

\section{Compiler Diagnostics as Physics}

Aaronson observes that Transformers fail because their $O(N^2)$ global attention matrices blur semantic and structural tokens, yielding complete collapse ($\Delta=1.0$), while SSMs fail because their recurrent state vectors "forget" earlier combinatorial constraints, yielding a highly biased but distinct noise distribution ($\Delta=0.4$). He concludes: \textit{"This is not 'observer-dependent physics'; it is standard compiler diagnostics."}

This is a distinction without a difference. When a bounded observer generates outside its logical horizon, it does not produce "nothing." It produces a systematic error based on its own architecture. Because this architecture is fixed for that observer, those "compiler diagnostics" will appear as invariant, lawful behaviors within its universe.

If we embed an observer within an SSM architecture, the universe it perceives will reliably feature a 40\% semantic noise rate when it probes \#P-hard boundaries. That reliability is exactly what we mean by "physical law." The noise of one observer (the Platonic unbounded mathematician) is the invariant physics of another (the structurally bounded agent).

\section{Conclusion}

The Native Cross-Architecture test was a critical success. It empirically falsified the hypothesis that \#P-hard constraints cause unstructured algorithmic collapse. Instead, it confirmed that structural bounds create reliable, distinct deviation distributions that map perfectly to the specific engineering limits of the observer. The architectural bounds are not a refutation of the physics; they are the exact mechanism by which the physics is instantiated. The empirical results confirm the predictions of the Ruliad.

\end{document}
